To project three-dimensional clustering information into the two-dimensional angular observables utilized by Stage-IV weak lensing and clustering likelihoods, your MCMC engine must evaluate line-of-sight integrals of the 3D matter power spectrum, weighted by adaptive geometric kernels.
Under the first-order Limber approximation, a generic angular cross-power spectrum $C_\ell^{XY}$ between two tomographic tracers $X$ and $Y$ (where $X, Y \in \{\gamma \text{ (shear)}, \delta_g \text{ (galaxy position)}, \kappa \text{ (CMB convergence)}\}$) at angular multipole $\ell$ is factually formulated as:
$$C_\ell^{XY} = \int_{0}^{\chi_{\rm *}} \frac{W_X(\chi) W_Y(\chi)}{\chi^2} P_{\delta \delta} \left( k = \frac{\ell + 1/2}{\chi}, z(\chi) \right) d\chi$$ 
Where:

* $\chi$ is the comoving distance acting as the line-of-sight integration variable.
* $\chi_{\rm *}$ is the comoving distance to the horizon (or surface of last scattering).
* $P_{\delta \delta}(k, z)$ is the non-linear 3D matter power spectrum evaluated at wavenumber $k = (\ell + 1/2)/\chi$.
* $W_X(\chi)$ and $W_Y(\chi)$ are the continuous, adaptive geometric weight kernels specific to each probe channel.

------------------------------
## 1. Mathematical Formulation of the Truncated Kernels
The exact mathematical structures of these weight kernels depend on how the physical geometry interacts with the distribution of the tracers along the line of sight:
## Cosmic Shear / Weak Lensing Kernel ($W_\gamma^i$)
For a given tomographic source bin $i$, the efficiency of gravitational lensing relies on the integrated mass distribution between the observer and the background source galaxies:
$$W_\gamma^i(\chi) = \frac{3 \Omega_m H_0^2}{2 c^2} \frac{\chi}{a(\chi)} \int_{\chi}^{\chi_{\rm *}} n_i(\chi') \frac{\chi' - \chi}{\chi'} d\chi'$$ 
Where $n_i(\chi')d\chi' = p_i(z)dz$ represents the normalized, selection-weighted photometric redshift distribution of galaxies assigned to tomographic bin $i$, and $a(\chi)$ is the scale factor.
## Galaxy Clustering / Position Kernel ($W_{\delta_g}^i$)
For a foreground lens galaxy sample, the kernel maps the localized spatial distribution of the tracer, modulated by its linear galaxy bias $b_i(z)$:
$$W_{\delta_g}^i(\chi) = b_i\left(z(\chi)\right) n_i\left(\chi\right) \frac{dz}{d\chi} = b_i\left(z(\chi)\right) n_i\left(\chi\right) H\left(z(\chi)\right)$$ 
## CMB Lensing Convergence Kernel ($W_\kappa$)
Because the Cosmic Microwave Background acts as a single, clean source plane located at the edge of the observable universe ($\chi_{\rm CMB} \approx 14,000 \text{ Mpc}$ at $z \sim 1100$), its geometric kernel does not use tomographic bins. Instead, it spans a broad, single-envelope kernel:
$$W_\kappa(\chi) = \frac{3 \Omega_m H_0^2}{2 c^2} \frac{\chi}{a(\chi)} \left( \frac{\chi_{\rm CMB} - \chi}{\chi_{\rm CMB}} \right)$$ 
------------------------------
## 2. Physical Evolution Across the 3x2pt Joint Matrix
The DES Y6 and KiDS-Legacy analyses build their entire multivariate Gaussian likelihood vectors by cross-correlating these exact kernels across different lines of sight. This is designated as the 3x2pt (three-times-two-point) configuration:

                    3x2pt OBSERVATIONAL PILLARS
                    
  [ W_γ × W_γ ]  ===> Cosmic Shear            ===> Pure mass-to-mass mapping
  [ W_δ_g × W_γ ] ===> Galaxy-Galaxy Lensing   ===> Lens-source cross-correlations
  [ W_δ_g × W_δ_g ]===> Galaxy Clustering      ===> Spatial autocorrelation of lenses


* Cosmic Shear ($C_\ell^{\gamma_i \gamma_j}$): Correlates two lensing kernels. Because the integral combines two broad background kernels ($\int_\chi^{\chi_*} n(\chi')d\chi'$), cosmic shear is highly sensitive to the smooth, integrated growth of dark matter structure across broad redshift windows.
* Galaxy-Galaxy Lensing ($C_\ell^{\delta_{g,i} \gamma_j}$): Correlates a sharp, localized galaxy position kernel ($W_{\delta_g}^i$) with a broad background lensing kernel ($W_\gamma^j$). This breaks the degeneracy between the true matter density and the unknown galaxy bias $b(z)$ by measuring the dark matter profile directly wrapped around the lens galaxies.
* Galaxy Clustering ($C_\ell^{\delta_{g,i} \delta_{g,j}}$): Correlates two identical galaxy position kernels. This provides a direct measure of the spatial power spectrum shape within narrow, localized spatial shells.

------------------------------
## 3. Adaptive Shift Implementation in the MCMC Chain
When checking a non-standard cosmological model, these line-of-sight integrals cannot be treated as static templates. The MCMC sampler must compute them adaptively at every single parameter coordinate step:

  [ Sampler Coordinate Step ] 
             │
             ▼
  [ Solve H(z), χ(z), and rd via Boltzmann ] 
             │
             ▼
  [ Adaptively deform W_X(χ) and W_Y(χ) integration grids ] 
             │
             ▼
  [ Shift non-linear scale matching: k = (ℓ + 1/2) / χ ] 
             │
             ▼
  [ Compute C_ℓ theoretical vector ] ---> [ Evaluate Matrix χ² Residual ]


   1. Geometric Scaling: If your model shifts $H_0$ or introduces dynamic dark energy ($w_0 w_a$), the comoving distance relation $\chi(z) = \int c dz/H(z)$ deforms. This dynamically shifts the peaks of the geometric kernels and alters the exact physical scale $k$ probed at a given angular multipole $\ell$.
   2. Growth Distortion: Modifications to gravity or early-universe parameters dynamically reshape the underlying 3D power spectrum $P_{\delta \delta}(k, z)$ over time. If your model suppresses structure growth at late times to resolve the $S_8$ tension, that reduction is integrated through the line-of-sight kernels, altering the final amplitude of the theoretical $C_\ell$ vector.
